% Generated draft: requires mathematical and source review.
\documentclass[11pt]{article}
\usepackage[margin=25mm]{geometry}
\usepackage{fontspec}
\setmainfont{Latin Modern Roman}
\usepackage{amsmath,amssymb,mathtools,unicode-math}
\setmathfont{Latin Modern Math}
\usepackage{enumitem,needspace,xurl,hyperref}
\setlength{\parindent}{0pt}
\setlength{\parskip}{5pt}
\setlength{\emergencystretch}{3em}
\newcommand{\N}{\mathbb{N}}
\newcommand{\Z}{\mathbb{Z}}
\newcommand{\Q}{\mathbb{Q}}
\newcommand{\R}{\mathbb{R}}
\newcommand{\C}{\mathbb{C}}
\newcommand{\F}{\mathbb{F}}
\newcommand{\A}{\mathbb{A}}
\newcommand{\PP}{\mathbb{P}}
\newcommand{\E}{\mathbb{E}}
\newcommand{\eps}{\varepsilon}
\newcommand{\dd}{\,\mathrm{d}}
\begin{document}
\section*{UnsolvedMath ID: 20002447}
\textbf{Source problem code:} catalogue entry
\subsection{Definitions and mathematical statement}
The source question is ``What is the stability of the KP-I soliton solution with periodic structure?'' The supplied verification places it under M. Haragus, ``Stability for KP-I profiles with periodic structure,'' and diagnoses a record-boundary error: the subsequent J. Albert Benjamin--Ono setup is a separate entry. That appended setup writes $u_t+uu_x-Ku_x=0$, $\widehat{Ku}(k)=|k|\widehat u(k)$ and $(cv+Kv-\phi v)_z=\lambda v$; its profile is rendered corruptly as ``41+ z2.'' These data do not define the KP-I wave and are not used as its equation.

The aid identifies the intended family as the $x$-localized, $y$-periodic Zaitsev traveling waves on $\R\times(\R/L\Z)$, $L>0$. The original extract supplies neither a KP-I normalization nor the family formula. For reference, in the common normalization
\[
 u_t+uu_x+u_{xxx}-\partial_x^{-1}u_{yy}=0,
\]
the natural energy norm is
\[
 \|u\|_{\mathcal E}^2=\|u\|_2^2+\|u_x\|_2^2+\|\partial_x^{-1}u_y\|_2^2.
\]
The inverse derivative is defined only for functions for which the Fourier multiplier expression is meaningful. The normalization and precise perturbation space must be checked against the cited results before using the numerical speeds below. Orbital stability in a specified norm means that every sufficiently small perturbation remains close for all forward times of existence to the spatial-translation orbit of the traveling wave; nonlinear global existence is part of the relevant stability theorem, not a consequence of this definition alone.

In the normalization of the aid's formulas, put $\kappa=2\pi/L$. Its reported local branch result is: for every $L>0$ there exists $a_*(L)>0$ such that for $0<|a|<a_*(L)$ the Zaitsev wave is orbitally stable and its speed satisfies
\[
 c_Z(a,L)-c_{\mathrm{crit}}(L)
 =\frac{\kappa a^2(2+a^2)}{\sqrt3(1-a^2)}>0,
 \qquad c_{\mathrm{crit}}(L)=\frac{8\pi}{\sqrt3 L},
\]
with the local parameter restricted also by $|a|<1$. The same cylinder has a $y$-independent KdV line soliton of speed $c_Z(a,L)$, and the aid reports it orbitally unstable because this speed lies above the line-soliton threshold. Thus the speed excess is relative to the critical line speed: the stable patterned wave and unstable line wave themselves have exactly the same speed.

The aid derives this corollary from Yamazaki's local Zaitsev orbital-stability theorem, the Rousset--Tzvetkov line-soliton threshold, and KP scaling. It advertises explicit arbitrary-period bifurcation-kernel modes without giving their formulas here. It does not settle stability along the entire Zaitsev branch toward the planar lump, which it describes as stable; the full-branch question remains unresolved in its dated assessment.
\subsection{Short English statement}
The KP-I question has been mixed with a following Benjamin--Ono entry. The supplied progress concerns only a small Zaitsev branch: at each fixed transverse period a stable patterned wave and an unstable line soliton can have the same speed. It does not establish stability of the whole branch.
\subsection{Sources}
\begin{itemize}
\item Local UnsolvedMath registry, permanent ID 20002447, source code catalogue entry; created and updated 14 August 2026. Canonical \texttt{aim-pdes-notes.json}, zero-based \texttt{notes[58]}, tag ``open problem,'' item 11.
\item AIM workshop ``Stability criteria for multi-dimensional waves and patterns,'' M. Haragus heading as reported in the supplied verification: \url{https://aimath.org/WWN/multidimwaves/multidimwaves.pdf}. The J. Albert Benjamin--Ono heading and setup belong to the following entry. Original remarks and literature field are empty.
\item Research attempt 1: \texttt{partial\_result}, problem status at run \texttt{partially\_solved}, rigor label \texttt{proved}, novelty confidence low. It names Yamazaki and Rousset--Tzvetkov, without complete paper titles, dates, theorem statements, or bibliographic links. No full authoritative citation for the local stability, threshold, or lump assessment is supplied.
\item Ulam AI / UnsolvedMath Contributors, CC BY 4.0. Machine-generated research aid, not peer reviewed; independently verify before citation.
\end{itemize}
\end{document}
